% ID: SEL-D0003-OB-E03
% ESTADO: cotejada. Fuente y fragmento cotejados de forma independiente; salida compilada revisada.
% Fuente: Selectividad/sel_2019_matematicas_aplicadas/Suplente Junio/Exámenes Matemáticas Aplicadas CCSS - A y B.pdf
% SHA-256: a59474e6895bd41034da53b7060969cbe8cbf7f8272f39fbf16821c95f06c84b
% Página: 2; opción B; ejercicio 3.
% Transcripción de primera pasada basada en la imagen de la página original.
% Se conserva la diferencia A-B y la barra de P(A/B^c) de la fuente.
% Sin solución ni modificación de datos o condiciones.

Sean \(A\) y \(B\) dos sucesos de un experimento aleatorio dado. Se sabe que \(P(A)=0.5\), \(P(A\cup B)=0.75\) y \(P(A-B)=0.3\).
\begin{enumerate}[label=\alph*),leftmargin=2.4em]
  \item \textbf{(0.5 puntos)} Calcule \(P(A\cap B)\).
  \item \textbf{(1 punto)} Calcule \(P(A/B^c)\).
  \item \textbf{(1 punto)} \textquestiondown{}Son independientes los sucesos \(A\) y \(B\)? \textquestiondown{}Son los sucesos \(A\) y \(B\) incompatibles?
\end{enumerate}
